\documentclass[10pt,a4paper]{amsart}
\usepackage[margin=19mm]{geometry}
\usepackage{amsmath,amssymb,mathtools}
\usepackage[scr=rsfs,cal=euler]{mathalfa}
\usepackage{xfrac}
\usepackage[unicode,hidelinks,bookmarks=false]{hyperref}
\newcommand{\ie}{i.e.\ }
\newcommand{\cf}{cf.\ }
\newcommand{\resp}{respectively\ }
\newcommand{\Subsection}{\subsection}
\providecommand{\nobreakdash}{\nobreak\hbox{-}}
\setlength{\emergencystretch}{2em}
\multlinegap0pt
\usepackage{amssymb}
\usepackage[shortlabels]{enumitem}
\usepackage{mathtools}
\usepackage{xcolor}
\newtheorem{lemma}{Lemma}
\newcommand{\R}{\mathbb{R}}
\def \d {{\rm{d}}}
\title{Quantitative observability for one-dimensional Schr\"odinger equations with potentials}
\author{Pei Su, Chenmin Sun, Xu Yuan}
\date{Source: arXiv:2309.00963v2 (CC BY 4.0)}
\begin{document}
\maketitle

\noindent\emph{Source and licence.} Quantitative observability for one-dimensional Schr\"odinger equations with potentials. Version arXiv:2309.00963v2 (CC BY 4.0), distributed by arXiv under the Creative Commons Attribution 4.0 International licence, http://creativecommons.org/licenses/by/4.0/, as recorded in the arXiv licence metadata for this exact version and shown on the abstract page https://arxiv.org/abs/2309.00963v2. Full licence notice: \texttt{licenses/arxiv-2309.00963-CC-BY-4.0.txt}.\par\smallskip

\noindent\textit{Editorial scope.} Counted unit: the article's lemma le:technical (source0.tex lines 637--642). The article's complete original proof of it is delivered with it (source0.tex lines 648--733, one \texttt{proof} environment of 86 lines, which the article places away from the statement). No mathematics is written, repaired or re-derived: the statement and the proof are the article's own lines, in the article's own notation, and no retained line is substituted or re-flowed.  No further source-local context is needed: every cross-reference of every retained line resolves inside the two delivered spans.  Nothing was omitted from any retained span, so the package carries no omission record.  No retained line contains a citation, so the wrapper carries no bibliography and no bibliography entry.  No retained line contains a figure environment, an image-inclusion command or drawing code, so no image is delivered and the package declares no asset dependency.  Editorial formatting only, in addition: an AMS \texttt{amsart} preamble that restates the article's own declarations and macros used by the retained lines, namely the environments \texttt{lemma} on separate counters, together with the standard packages those lines need.  The retained environments print in this document as Lemma 1 in order of appearance, whereas the article prints the same items with the numbering of its own full document.  The complete native source of the article as distributed in the arXiv e-print for this exact version is bundled unchanged as \texttt{sources/146455/source0.tex}.
\begin{lemma}\label{le:technical}
    Let $0<\zeta<1$. Then there exists a constant $c=c(\zeta)>0$, depending only on $\zeta$, such that for any measurable set $\omega\subset (0,1)$ with $\zeta\le |\omega|\le 1$ and any $\lambda>1$, we have 
    \begin{equation}\label{est:A}
       \inf_{x_{0}\in \R} \int_{\omega}\cos^{2}(\lambda(x-x_{0}))\d x\ge c(\zeta).
    \end{equation}
\end{lemma}

\begin{proof}[Proof of Lemma~\ref{le:technical}]
Using the structure of a measurable set in $\R$, for any $0<\varepsilon<1$, there exists a finite union of disjoint open intervals 
\begin{equation*}
    U=\bigcup\limits_{n=1}^{N}I_{n}\quad \mbox{with}\quad 
    I_{n}=(a_{n},b_{n})\subset (0,1) \  \ \mbox{for any}\ n\in \left\{1,\dots,N\right\},
    \end{equation*}
    such that $|\omega \setminus U|+|U\setminus \omega |<\varepsilon$ and $|U|\ge \zeta-\varepsilon$.
    It follows that 
    \begin{equation}\label{est:cos2}
    \begin{aligned}
       \int_{\omega}\cos^{2}(\lambda(x-x_{0}))\d x
       &\ge 
       \sum_{n=1}^{N}\int_{a_{n}}^{b_{n}}\cos^{2}(\lambda(x-x_{0}))\d x-\varepsilon.
       \end{aligned}
        \end{equation}
By an elementary computation, on any finite interval $I_{n}=(a_{n},b_{n})$, we have 
\begin{equation}\label{est:ajbj}
\begin{aligned}
    &\int_{a_{n}}^{b_{n}}\cos^{2}(\lambda(x-x_{0}))\d x\\
    &=\frac{1}{2}\int_{a_{n}}^{b_{n}}(1+\cos (2\lambda(x-x_{0})))\d x\\
    &=\frac{1}{2}\left((b_{n}-a_{n})+\frac{1}{\lambda}\sin \left(\lambda(b_{n}-a_{n})\right)\cos \left(\lambda(b_{n}+a_{n}-2x_{0})\right)\right).
    \end{aligned}
    \end{equation}
Let $0<\delta\ll 1$ be a small constant to be chosen later. For any $m\in \mathbb{Z}$, we set 
\begin{equation*}
    J_{m,\delta}=\left(x_{0}+\frac{m \pi }{2\lambda}-\frac{\delta}{2\lambda},x_{0}+\frac{m \pi }{2\lambda}+\frac{\delta}{2\lambda}\right)\cap (0,1)\ \ \mbox{and}\  \ S=\bigcup\limits_{m\in \mathbb{Z}}J_{m,\delta}.
    \end{equation*}
Note that, $\left\{J_{m,\delta}\right\}_{m\in \mathbb{Z}}$ are disjoint sets.
For further reference, we consider
\begin{equation*}
    \mathcal{A}=\left\{m\in \mathbb{Z}:J_{m,\delta}\ne \emptyset\right\},\quad \mbox{and thus we have}\quad  \# \mathcal{A}\le 3\lambda.
\end{equation*}
We split the index set $\left\{1,\dots,N\right\}$ to the following three cases according to the size and position of the finite interval $I_{n}$ and then establish estimate for each case.

\smallskip
\textbf{Case 1.} Let $(b_{n}-a_{n})\ge \frac{\delta}{\lambda}$. From the fact that $\frac{\sin x}{x}$ is decreasing on $[\delta,\pi)$, for $0<\delta\ll 1$, 
\begin{equation*}
\sup_{[\delta,\infty)}\left|\frac{\sin x}{x}\right|\le \frac{\sin \delta}{\delta} \Longrightarrow \frac{1}{\lambda}\left|\sin \left(\lambda(b_{n}-a_{n})\right)\right|\le \frac{\sin \delta}{\delta}(b_{n}-a_{n}).
\end{equation*}
Based on the above estimate and~\eqref{est:ajbj}, we obtain
\begin{equation}\label{est:case1}
\int_{a_{n}}^{b_{n}}\cos^{2}(\lambda(x-x_{0}))\d x
\ge \frac{1}{2}\left(1-\frac{\sin \delta}{\delta}\right)(b_{n}-a_{n}).
\end{equation}

\textbf{Case 2.} Let $0<(b_{n}-a_{n})< \frac{\delta}{\lambda}$ with $S\cap I_{n}= \emptyset$. First, from the definition of $S$ and $S\cap I_{n}=\emptyset$, there exists $m\in \mathbb{Z}$ such that 
\begin{equation*}
    I_{n}\subset \left(x_{0}+\frac{m \pi }{2\lambda}+\frac{\delta}{2\lambda},x_{0}+\frac{ (m+1)\pi}{2\lambda}-\frac{\delta}{2\lambda}\right),
    \end{equation*}
    and thus, from $\sin^{2}x+\cos^{2}x=1$, we find
    \begin{equation*}
    \begin{aligned}
   \mbox{dist}(\lambda(b_{n}+a_{n}-2x_{0}),\pi\mathbb{Z})\ge \delta\Longrightarrow
    \left|\cos (\lambda(b_{n}+a_{n}-2x_{0})\right|\le \sqrt{1-\sin^{2}\delta}.
    \end{aligned}
    \end{equation*}
    Therefore, using again~\eqref{est:ajbj} and $|\sin x|\le |x|$, we obtain
    \begin{equation}\label{est:case2}
        \int_{a_{n}}^{b_{n}}\cos^{2}(\lambda(x-x_{0}))\d x\ge
        \frac{1}{2}\left(1-\sqrt{1-\sin^{2}\delta}\right)(b_{n}-a_{n}).
        \end{equation}

\textbf{Case 3.} We now consider the last case, that is, the case of $0<(b_{n}-a_{n})< \frac{\delta}{\lambda}$ with $S\cap I_{n}\ne  \emptyset$. To simplify notation, we denote 
\begin{equation*}
    \mathcal{B}=\left\{n\in\left\{1,\dots,N\right\}:0<(b_{n}-a_{n})< \frac{\delta}{\lambda}\ \ \mbox{with} \ \ S\cap I_{n}\ne  \emptyset\right\}.
\end{equation*}

Note that, for any $n\in \mathcal{B}$, there exists $m\in\mathcal{A}$ such that 
\begin{equation*}
I_{n}\subset {J}_{m,3\delta}\quad \mbox{where}\ \ 
{J}_{m,3\delta}=\left(x_{0}+\frac{m \pi }{2\lambda}-\frac{3\delta}{2\lambda},x_{0}+\frac{m \pi }{2\lambda}+\frac{3\delta}{2\lambda}\right)\cap(0,1).
\end{equation*}
Next, from $0<\delta\ll 1$, for any $(m,m')\in \mathbb{Z}^{2}$ with $m\ne m'$, we find ${J}_{m,3\delta}\cap
{J}_{m',3\delta}=\emptyset$.
Therefore, using the fact that $\left\{I_{n}\right\}_{n=1}^{N}$ are disjoint intervals and $\# \mathcal{A}\le 3\lambda$, we obtain
    \begin{equation}\label{est:case3}
   \sum_{n\in \mathcal{B}}(b_{n}-a_{n})=\bigg|\bigcup\limits_{n\in \mathcal{B}}I_{n}\bigg|\le \bigg|\bigcup\limits_{m\in \mathcal{A}}{J}_{m,3\delta}\bigg|\le \frac{3\delta}{\lambda}\#\mathcal{A}\le  9\delta.
   \end{equation}
   

   Combining~\eqref{est:cos2},~\eqref{est:case1},~\eqref{est:case2},~\eqref{est:case3} with $|U|\ge \zeta-\varepsilon$, we conclude that 
   \begin{equation*}
    \inf_{x_{0}\in \R}\int_{\omega}\cos^{2}(\lambda(x-x_{0}))\d x\ge \frac{1}{2}\left(1-\max\left(\frac{\sin \delta}{\delta},\sqrt{1-\sin^{2}\delta}\right)\right)(\zeta-\varepsilon-9\delta)-\varepsilon.
   \end{equation*}
   We see that~\eqref{est:A} follows from the above estimate fo $\varepsilon$ and $\delta$ small enough.
\end{proof}

\end{document}
